\documentclass[12pt]{article}
\usepackage[utf8]{inputenc}
\usepackage[T1]{fontenc}
\usepackage{hyperref}
\title{Vanuma Language Resources for K-12 Curriculum Development}
\author{Compiled by Claude Sonnet 5 (claude-sonnet-5) for Andr\'as Kornai}
\date{}
\begin{document}
\maketitle

\section{Introduction}
Vanuma (Bvanuma, South Nyali; ISO 639-3 \texttt{vau}, Glottocode vanu1242) is a Bantu language (Guthrie D.331) of Ituri Province, Democratic Republic of the Congo. About 6,700 speakers are cited, from a 1993 count. Very little published material was found.

\section{Primary Educational Resources}
No primers, textbooks or school programmes were found in the search results.

\section{Linguistic and Reference Materials}
\begin{itemize}
\item Wikipedia overview: \url{https://en.wikipedia.org/wiki/Vanuma_language}
\item ASJP database entry (word list data): \url{https://asjp.clld.org/languages/VANUMA}
\item Conference presentations by Douglas W.\ Boone (SIL), which the search snippet describes as including work on Vanuma endangerment, interviews with Vanuma children in 2014 and a 2021 literacy study (contents not inspected): \url{https://sil.academia.edu/DouglasWBoone}
\end{itemize}

\section{Religious and Community Resources}
\begin{itemize}
\item Joshua Project profile of the Vanuma: \url{https://joshuaproject.net/people_groups/15773}
\end{itemize}

\section{Digital and Audio Resources}
\begin{itemize}
\item Wikitongues language page (listing only; no recordings confirmed): \url{https://wikitongues.org/languages/vau/}
\end{itemize}

\section{Curriculum Development Context}
The language is reported as not taught in schools. The search found no dictionary, grammar, primer or orthography. The Boone literacy work is the most relevant lead and should be followed up directly. Any curriculum would require substantial new documentation with the community.
\end{document}